\subsection{A fourteen-dimensional extension \texorpdfstring{$t$}{t}-module}
\label{sec:large-example}
\setcounter{MaxMatrixCols}{14}

Let $L=\Fq(\theta)$, with $\theta$ transcendental over $\Fq$,
and consider
\begin{equation}
 \Phi_t=
 \begin{pmatrix}
  \theta+\tau^4&\tau\\
  0&\theta+\tau^3
 \end{pmatrix},
 \qquad
 \Psi_t=
 \begin{pmatrix}
  \theta+\theta^2\tau&0\\
  1+\tau&\theta+\theta\tau
 \end{pmatrix}.
 \label{eq:large-presentations}
\end{equation}
The coefficient of $\tau^4$ in $\Phi_t$ has rank one.
The nonzero positive coefficient blocks are
\[
 A_1=\begin{pmatrix}0&1\\0&0\end{pmatrix},\qquad
 A_3=\begin{pmatrix}0&0\\0&1\end{pmatrix},\qquad
 A_4=\begin{pmatrix}1&0\\0&0\end{pmatrix}.
\]
Here $A_2=0$. The suffix ranks are $s_4=1$ and $s_3=2$,
so Algorithm~\ref{alg:preprocess} retains $[A_3\mid A_4]$
and returns
\[
 \mathrm{Flag}=((3,2),(4,1)),\qquad
 \mathrm{Required}=(0,0,1,1).
\]
The candidate lists contain only $c_4^1=e_1$ and
$c_3^2=e_2$. Their compatible-order graph has no edges,
so Algorithm~\ref{alg:order} returns
\begin{equation}
 \sigma=(1,2),\qquad
 \mathbb B=((1,4,e_1),(2,3,e_2)),\qquad C=I_2.
 \label{eq:large-compatible-basis}
\end{equation}
The selected vectors form the compatible basis
$(c_4^1,c_3^2)=(e_1,e_2)$.
Indeed, reading the retained blocks in this order gives
\[
 (c_4^1,c_4^2,c_3^1,c_3^2)=(e_1,0,0,e_2),
\]
with successive ranks $1,1,1,2$. Thus
\[
 (j_1,j_2)=(4,3),\qquad N=3,\qquad J=7.
\]
Consequently, $\Phi$ is of type $\mathrm{FJ}_{(1,2),3}$,
and the degree gap holds because $\deg_\tau\Psi_t=1<3$.

The nilpotent parts are
\[
 N_\Phi=0,\qquad
 N_\Psi=\begin{pmatrix}0&0\\1&0\end{pmatrix}.
\]
Since $C=I_2$, the intertwining condition is
$N_\Psi E_{i,s}=0$. It holds precisely for $(i,s)=(2,1)$
and $(2,2)$. Hence Algorithm~\ref{alg:ext0-gapart} gives
\begin{equation}
 \varepsilon=\begin{pmatrix}1&1\\0&0\end{pmatrix},
 \qquad h=2,\qquad M=14,\qquad M_0=12.
 \label{eq:large-dimensions}
\end{equation}

The full normal representatives satisfy
$\deg_\tau F_{i,1}<4$ and $\deg_\tau F_{i,2}<3$.
We use the ordered basis
\[
\begin{aligned}
 \EE=\big(&E_{11},E_{11}\tau,E_{11}\tau^2,E_{11}\tau^3,
            E_{12},E_{12}\tau,E_{12}\tau^2,\\
          &E_{21},E_{21}\tau,E_{21}\tau^2,E_{21}\tau^3,
            E_{22},E_{22}\tau,E_{22}\tau^2\big).
\end{aligned}
\]
For $\Ext_0(\Phi,\Psi)$, the ordered family is
\[
\begin{aligned}
 \mathbf E_0=\big(&E_{11}\tau,E_{11}\tau^2,E_{11}\tau^3,
                   E_{11}\tau^4,
                   E_{12}\tau,E_{12}\tau^2,E_{12}\tau^3,\\
                 &E_{21}\tau,E_{21}\tau^2,E_{21}\tau^3,
                   E_{22}\tau,E_{22}\tau^2\big).
\end{aligned}
\]
The additional degrees four and three in the first row
come from $\varepsilon_{1,1}=\varepsilon_{1,2}=1$.

To keep the displayed matrices readable, put
\[
 a=\theta^2,\qquad b=\theta,\qquad
 \Delta=\theta^{(1)}-\theta,
\]
and write
\begin{equation}
\begin{aligned}
 \alpha&=a a^{(1)}\tau^2,
 &\beta&=b b^{(1)}\tau^2,
 &\gamma&=(b+a^{(1)})\tau^2,\\
 \xi&=-a\Delta\tau,
 &\omega&=(a+b-\Delta)\tau.
\end{aligned}
 \label{eq:large-entry-notation}
\end{equation}
Here $u^{(1)}=u^q$, and coefficients are written on the
left of $\tau$. Algorithms~\ref{alg:reduce} and
\ref{alg:extmatrix} give the full action matrix
\begingroup
\setlength{\arraycolsep}{2pt}
\begin{equation}
 \Pi_t=
\begin{pmatrix}
  \theta & 0 & 0 & 0 & 0 & 0 & 0 & 0 & 0 & 0 & 0 & 0 & 0 & 0\\
  a\tau & \theta & 0 & \alpha & 0 & 0 & 0 & 0 & 0 & 0 & 0 & 0 & 0 & 0\\
  0 & a\tau & \theta & 0 & 0 & 0 & 0 & 0 & 0 & 0 & 0 & 0 & 0 & 0\\
  0 & 0 & a\tau & \theta & 0 & 0 & 0 & 0 & 0 & 0 & 0 & 0 & 0 & 0\\
  0 & 0 & 0 & 0 & \theta & 0 & 0 & 0 & 0 & 0 & 0 & 0 & 0 & 0\\
  0 & 0 & 0 & -a\tau & a\tau & \theta & \alpha & 0 & 0 & 0 & 0 & 0 & 0 & 0\\
  0 & 0 & 0 & 0 & 0 & a\tau & \theta & 0 & 0 & 0 & 0 & 0 & 0 & 0\\
  1 & 0 & 0 & a\tau & 0 & 0 & 0 & \theta & 0 & 0 & 0 & 0 & 0 & 0\\
  \tau & 1 & 0 & \gamma & 0 & 0 & 0 & b\tau & \theta & 0 & \beta & 0 & 0 & 0\\
  0 & \tau & 1 & 0 & 0 & 0 & 0 & 0 & b\tau & \theta & 0 & 0 & 0 & 0\\
  0 & 0 & \tau & 1 & 0 & 0 & 0 & 0 & 0 & b\tau & \theta & 0 & 0 & 0\\
  0 & 0 & 0 & 0 & 1 & 0 & a\tau & 0 & 0 & 0 & 0 & \theta & 0 & 0\\
  0 & 0 & 0 & -\tau & \tau & 1 & \gamma & 0 & 0 & 0 & -b\tau & b\tau & \theta & \beta\\
  0 & 0 & 0 & 0 & 0 & \tau & 1 & 0 & 0 & 0 & 0 & 0 & b\tau & \theta
 \end{pmatrix}
 \label{eq:large-full-action}
\end{equation}
\endgroup
Its differential is
\[
 \partial\Pi_t=
 \begin{pmatrix}\theta I_7&0\\I_7&\theta I_7\end{pmatrix}
 = (\theta I_2+N_\Psi)\otimes I_7.
\]

Algorithms~\ref{alg:ext0-reduce} and~\ref{alg:ext0-action}
give the action on $\Ext_0$ in the order $\mathbf E_0$:
\begingroup
\setlength{\arraycolsep}{2pt}
\begin{equation}
 \Pi_{0,t}=
\begin{pmatrix}
  \theta & 0 & 0 & \xi & 0 & 0 & 0 & 0 & 0 & 0 & 0 & 0\\
  a\tau & \theta & 0 & \alpha & 0 & 0 & 0 & 0 & 0 & 0 & 0 & 0\\
  0 & a\tau & \theta & 0 & 0 & 0 & 0 & 0 & 0 & 0 & 0 & 0\\
  0 & 0 & a\tau & \theta & 0 & 0 & 0 & 0 & 0 & 0 & 0 & 0\\
  0 & 0 & 0 & 0 & \theta & 0 & \xi & 0 & 0 & 0 & 0 & 0\\
  0 & 0 & 0 & -a\tau & a\tau & \theta & \alpha & 0 & 0 & 0 & 0 & 0\\
  0 & 0 & 0 & 0 & 0 & a\tau & \theta & 0 & 0 & 0 & 0 & 0\\
  1 & 0 & b\tau^2 & \omega & 0 & 0 & 0 & \theta & 0 & \beta & 0 & 0\\
  \tau & 1 & 0 & \gamma & 0 & 0 & 0 & b\tau & \theta & 0 & 0 & 0\\
  0 & \tau & 1 & 0 & 0 & 0 & 0 & 0 & b\tau & \theta & 0 & 0\\
  0 & 0 & -\tau & -1 & 1 & b\tau^2 & \omega & 0 & 0 & -b\tau & \theta & \beta\\
  0 & 0 & 0 & -\tau & \tau & 1 & \gamma & 0 & 0 & 0 & b\tau & \theta
 \end{pmatrix}
 \label{eq:large-ext0-action}
\end{equation}
\endgroup
The constant part of this matrix is $\theta I_{12}$ plus
a square-zero matrix. Thus both displayed actions satisfy
the Anderson condition, with nonzero nilpotent parts.

Only the fourth and seventh elements of $\mathbf E_0$
require reduction to obtain their full normal coordinates.
For a formal scalar $c$, these reductions are
\begin{align}
 \operatorname{red}(cE_{11}\tau^4)
   &=a c^qE_{11}\tau-cE_{12}\tau
       +cE_{21}+c^qE_{21}\tau,\\
 \operatorname{red}(cE_{12}\tau^3)
   &=a c^qE_{12}\tau+cE_{22}+c^qE_{22}\tau.
 \label{eq:large-boundary-reductions}
\end{align}
They follow by subtracting the inner biderivations
determined by $cE_{11}$ and $cE_{12}$, respectively.
The terms $c^q$ give entries containing $\tau$ in the
inclusion matrix. Algorithm~\ref{alg:ext0-transition}
therefore returns
\begingroup
\setlength{\arraycolsep}{2pt}
\begin{equation}
 T=
\begin{pmatrix}
  0 & 0 & 0 & 0 & 0 & 0 & 0 & 0 & 0 & 0 & 0 & 0\\
  1 & 0 & 0 & a\tau & 0 & 0 & 0 & 0 & 0 & 0 & 0 & 0\\
  0 & 1 & 0 & 0 & 0 & 0 & 0 & 0 & 0 & 0 & 0 & 0\\
  0 & 0 & 1 & 0 & 0 & 0 & 0 & 0 & 0 & 0 & 0 & 0\\
  0 & 0 & 0 & 0 & 0 & 0 & 0 & 0 & 0 & 0 & 0 & 0\\
  0 & 0 & 0 & -1 & 1 & 0 & a\tau & 0 & 0 & 0 & 0 & 0\\
  0 & 0 & 0 & 0 & 0 & 1 & 0 & 0 & 0 & 0 & 0 & 0\\
  0 & 0 & 0 & 1 & 0 & 0 & 0 & 0 & 0 & 0 & 0 & 0\\
  0 & 0 & 0 & \tau & 0 & 0 & 0 & 1 & 0 & 0 & 0 & 0\\
  0 & 0 & 0 & 0 & 0 & 0 & 0 & 0 & 1 & 0 & 0 & 0\\
  0 & 0 & 0 & 0 & 0 & 0 & 0 & 0 & 0 & 1 & 0 & 0\\
  0 & 0 & 0 & 0 & 0 & 0 & 1 & 0 & 0 & 0 & 0 & 0\\
  0 & 0 & 0 & 0 & 0 & 0 & \tau & 0 & 0 & 0 & 1 & 0\\
  0 & 0 & 0 & 0 & 0 & 0 & 0 & 0 & 0 & 0 & 0 & 1
 \end{pmatrix}
 \label{eq:large-inclusion}
\end{equation}
\endgroup

For an inner biderivation, the constant coefficient is
$-N_\Psi H_0$, which is supported in the second row.
The quotient map consequently records the constant
coefficients in the first row. In the coordinate order
$\EE$, it is
\begin{equation}
 P=
\begin{pmatrix}
  1 & 0 & 0 & 0 & 0 & 0 & 0 & 0 & 0 & 0 & 0 & 0 & 0 & 0\\
  0 & 0 & 0 & 0 & 1 & 0 & 0 & 0 & 0 & 0 & 0 & 0 & 0 & 0
 \end{pmatrix}
 \label{eq:large-projection}
\end{equation}
Direct Ore multiplication gives
\begin{equation}
 \Pi_tT=T\Pi_{0,t},\qquad
 P\Pi_t=\theta P,\qquad PT=0.
 \label{eq:large-intertwining}
\end{equation}

Exactness can also be checked directly in these coordinates.
If $z=(z_1,\ldots,z_{14})$ satisfies $Pz=0$, then
$z_1=z_5=0$ and its unique preimage under $T$ is
\[
\begin{aligned}
 (&z_2-a z_8^q,\ z_3,\ z_4,\ z_8,\
   z_6+z_8-a z_{12}^q,\ z_7,\ z_{12},\\
  &z_9-z_8^q,\ z_{10},\ z_{11},\
   z_{13}-z_{12}^q,\ z_{14}).
\end{aligned}
\]
This is a polynomial inverse on $\ker P$. The map $P$ is
surjective, so the full exact sequence in normal coordinates is
\begin{equation}
 0\longrightarrow(\Ga^{12},\Pi_{0,t})
 \xrightarrow{T}(\Ga^{14},\Pi_t)
 \xrightarrow{P}(\Ga^2,\theta I_2)
 \longrightarrow0.
 \label{eq:large-coordinate-sequence}
\end{equation}
Under the identifications with extension classes, this becomes
\begin{equation}
 0\longrightarrow\Ext_0(\Phi,\Psi)
 \xrightarrow{\iota}\Ext^1_\tau(\Phi,\Psi)
 \xrightarrow{p}(\Ga^2,\theta I_2)
 \longrightarrow0.
 \label{eq:large-extension-sequence}
\end{equation}

The search for this example has one candidate in each
list and tests one completed selection. Across all columns,
Algorithms~\ref{alg:extmatrix}, \ref{alg:ext0-action}, and
\ref{alg:ext0-transition} use six, ten, and two nonzero inner
corrections, respectively, compared with the uniform bounds
$14K(4)=84$, $12K_0(5)=96$, and $2K(4)=12$ obtained from
\eqref{eq:complexity-reduction-counts}. The two action matrices have
$\tau$-degree two, and the inclusion has $\tau$-degree one.
